\problemname{Perpetuum Mobile}
The year is 1902. Albert Einstein is working in the patent office in Bern.
Many patent proposals contain egregious errors; some even violate the law of conservation of energy.
To make matters worse, the majority of proposals make use of non-standard 
physical units that are not part of the metric system (or not even documented).
All proposals are of the following form:

\begin{itemize}
	\item Every patent proposal contains $n$ energy converters.
	\item Every converter has an unknown input energy unit associated with it.
	\item Some energy converters can be connected: If converter $a$ can be connected to 
	      converter $b$ such that one energy unit associated with $a$ is turned into 
	      $c$ input units for $b$, then this is indicated by an arc 
	      $a \overset{c}{\longrightarrow} b$ in the proposal. 
	      The output of $a$ can be used as input for $b$ 
	      if and only if such an arc from $a$ to $b$ exists.
\end{itemize}
Einstein would like to dismiss all those proposals out of hand where the energy 
converters can be chained up in a cycle such that more energy is fed back to a converter 
than is given to it as input, thereby violating the law of conservation of energy.

Einstein's assistants know that he is born for higher things than weeding out faulty 
patent proposals. Hence, they take care of the most difficult cases, while 
the proposals given to Einstein are of a rather restricted form: 
Every \textit{admissible} patent proposal given to Einstein 
does not allow for a cycle where the total product of arc weights exceeds $0.9$.
By contrast, every \textit{inadmissible} patent proposal given to Einstein contains a cycle 
where the the number of arcs constituting the cycle does not exceed the number of converters defined in the proposal, 
and the total product of arc weights is greater or equal to $1.1$. 

Could you help Einstein identify the inadmissible proposals?
\section*{Input}
The input consists of:
\begin{itemize}
	\item one line with two integers $n$ and $m$, where
		\begin{itemize}
			\item $n$ ($2\leq n \leq 800$) is the number of energy converters;
			\item $m$ ($0 \leq m \leq 4000$) is the number of arcs.
		\end{itemize}
	\item $m$ lines each containing three numbers $a_i$, $b_i$, and $c_i$, where
		\begin{itemize}
			\item $a_i$ and $b_i$ ($1 \leq a_i, b_i \leq n$) are integers identifying energy converters;
			\item $c_i$ ($0 < c_i \leq 5.0$) is a decimal number
    \end{itemize}
    indicating that the converter $a_i$ can be connected to the converter $b_i$ such that one input unit associated with $a_i$ is converted to $c_i$ units associated with $b_i$. The number $c_i$ may have up to $4$ decimal places.
\end{itemize}

%The input starts with a single line containing two space-separated 
%integers $n \, m$ $(2\leq n \leq 800, 0 \leq m \leq 800)$ where $n$ is the number of energy converters and $m$ is the 
%number of arcs. Then $m$ lines follow. The $i$th line is of the form 
%$a_i \, b_i \, c_i$ $(1 \leq a_i, b_i \leq n, 0.0001 \leq c_i \leq 5.0)$ where $a_i$ and $b_i$ are integers identifying energy converters and $c_i$ 
%is a decimal number, indicating that the converter $a_i$ can be connected to the 
%converter $b_i$ such that one input unit associated with $a_i$ 
%is converted to $c_i$ units associated with $b_i$. The number $c_i$ may have up to $4$ 
%decimal places.

\section*{Output}

Output a single line containing \texttt{inadmissible} if the proposal given to Einstein is inadmissible, \texttt{admissible} otherwise. 
